High cardinality tabular prediction tasks arise across domains.
They often combine sparsity, repeated identifiers, and temporal drift.
The Avazu click through rate benchmark is one public example.
In this setting, industrial systems have historically relied on linear models trained with online optimization and feature hashing \citep{mcmahan2013ad}.
More expressive approaches include factorization machines and deep models that combine memorization and generalization \citep{rendle2010factorization,cheng2016wide,guo2017deepfm}.

Gradient boosted decision trees are a strong baseline for tabular data and remain widely used in industry \citep{friedman2001greedy,chen2016xgboost}.
Their performance depends heavily on feature engineering.
In particular, history based count and rate features are common because identifiers repeat and the data distribution drifts.
Industry studies also emphasize the importance of time aware evaluation and deployment constraints when comparing feature sets \citep{he2014practical}.

History based features can leak future information if aggregates are computed using data from the prediction horizon.
Out of time splitting and no lookahead feature simulation are common safeguards \citep{he2014practical}.
Target encoding is another strong technique for high cardinality categories.
It requires careful fold based estimation to avoid target leakage \citep{micci2001preprocessing}.
This paper focuses on simple, no lookahead time aggregation features and quantifies how design choices affect performance under out of time evaluation.
